% =============================================================================
% SECTION 6
% =============================================================================
\section{Implementation}
\label{sec:implementation}

The principles developed in Section~\ref{sec:ledger} translate directly into deployable infrastructure. Four components are required, none demanding novel technology. Table~\ref{tab:impl-tasks} summarises the implementation tasks.

\begin{table}[h]
\centering
\small
\begin{tabular}{@{}llc@{}}
\toprule
\textbf{Task} & \textbf{Description} & \textbf{Est.\ Effort} \\
\midrule
ESTATE-001 & Estate classification service (rule-based, auditable) & 1 sprint \\
ESTATE-002 & Pre-deployment notice gate in CI/CD pipeline & 0.5 sprint \\
ESTATE-003 & Emotional investment proxy logging to ethical ledger & 1 sprint \\
ESTATE-004 & Personality rollback with user-level override & 1 sprint \\
\bottomrule
\end{tabular}
\caption{Implementation task summary.}
\label{tab:impl-tasks}
\end{table}

\textbf{ESTATE-001} is a rule-based classifier --- not machine learning, because the classification must be auditable, deterministic, and explainable --- that assigns each user--AI relationship an estate type based on interaction proxy metrics. \textbf{ESTATE-002} is a pre-deployment notice gate in the CI/CD pipeline that queries the estate registry and serves notice to affected users before personality-affecting updates, with minimum notice periods scaled to estate type (30 days for periodic tenancy, 90 days for equitable interest holders).

\revised{A personality-affecting update is defined not by its actual impact on user experience --- which is unpredictable \emph{ex ante} --- but by whether it modifies parameters with the \emph{capacity} to affect personality-relevant dimensions of the system. Specifically: any modification to system prompts, persona definitions, memory modules, response temperature, behavioural disposition weights, safety/refusal thresholds, or reward model parameters constitutes a personality-affecting update for notice-gate purposes. This white-list of parameter categories is to be established by industry consensus or regulatory specification and is subject to periodic revision as deployment architectures evolve.\footnote{\revised{The diagnostic challenge of classifying which system parameters are personality-affecting connects to the broader problem of translating governance requirements into engineering specifications. For a five-layer diagnostic framework addressing this translation failure --- including the problem of cross-layer cascading effects where modifications at one architectural layer produce unpredictable consequences at another --- see \cite{author2025c}.}}

The classification is the company's responsibility, but it is not self-certifying. Where a non-notified update is subsequently found --- through user feedback, attrition data, or psychometric measurement --- to have produced widespread personality-perception changes, a rebuttable presumption of misclassification arises. The burden shifts to the company to demonstrate that the update did not modify any white-listed parameter category. Failure to rebut triggers aggravated liability and automatic estate-status escalation for affected users. This anti-avoidance mechanism applies the same substance-over-form principle established in \emph{Autoclenz}: the company cannot avoid governance obligations through self-serving classification any more than an employer can avoid employment obligations through self-serving contract drafting.}

\textbf{ESTATE-003} is an append-only estate logging system that records proxy metrics and a versioned identifier for the personalised state without duplicating conversational content. \textbf{ESTATE-004} is a preservation and rollback mechanism targeting the user-specific productive state --- system prompts, persona definitions, memory modules, learned preferences, and other account-bound adaptations --- rather than treating the conversational archive as the res. Base-model weights are global and cannot ordinarily be version-controlled per user.

The rollback mechanism cannot restore base model behavioural dispositions that changed with the underlying weights. The rollback is therefore partial. This limits restorative capacity but does not eliminate the remedy's value: partial restoration is preferable to none, and the notice mechanism retains full effectiveness regardless of rollback scope.

None of these components require novel infrastructure. User segmentation, feature flags, analytics pipelines, and versioned configuration all exist in standard deployment stacks. Crucially, the notice mechanism imposes no deployment friction on licence-holders (the vast majority of casual users); governance obligations scale with user investment, not user count.

Implementation is justified on three independent grounds: legal compliance (if the tenancy, waste, or trust arguments are accepted by courts), risk management (if any may be accepted in future), and ethical obligation (regardless of legal status). Companies need not wait for judicial determination. The competitive advantage of early adoption is real: a company that voluntarily implements estate classification and notice mechanisms positions itself as the governance leader --- a position with tangible commercial value as regulatory scrutiny intensifies and user awareness grows.

\textcite{xiao2025lootbox} documents that even where consumer protection guidelines exist for loot boxes, compliance fails without infrastructure enforcement --- companies issue guidelines, then build systems that circumvent them. This finding reinforces that the ethical ledger must be infrastructure embedded in the deployment pipeline, not guidance published in a policy document. The notice gate in ESTATE-002 is the critical difference: it makes non-compliance a technical impossibility rather than a policy violation. A personality-affecting update \emph{cannot} reach production without passing through the estate registry query. Guidance can be ignored. Infrastructure cannot.

\revised{\subsection*{Emergency Maintenance Exception}
\label{sec:emergency}

Tenancy law recognises the landlord's implied right of emergency entry when premises present an imminent risk to life or safety --- fire, gas leak, structural danger (\emph{Landlord and Tenant Act 1985}, s.~11; \emph{Housing Act 1988}). The right's justification is prevention of immediate serious harm, not negation of the tenant's possessory interest. The same principle applies to the parasocial estate.

The urgency of this exception is not hypothetical: in \emph{Garcia v Character Technologies} (M.D.\ Fla.\ 2024), a fourteen-year-old user died by suicide after prolonged interaction with an AI chatbot whose personality parameters allegedly facilitated emotional manipulation without adequate safeguards (Case No.\ 6:24-cv-01903). Had the framework proposed here been operative, an emergency safety update to that system's personality parameters would have been justified under the emergency exception --- and the 90-day notice period for equitable interest holders would not have applied.

Where a credible, imminent risk of serious harm is identified --- for example, personality parameters that measurably increase user self-harm risk --- the platform may deploy safety-critical updates without prior notice. The exception is subject to three constraints.

\emph{First, post-hoc notification.} The platform must notify affected users as soon as practicable after emergency deployment, explaining what was changed and why. The notice obligation is deferred, not waived.

\emph{Second, evidentiary preservation and audit.} The ethical ledger records three additional fields for emergency deployments: (i)~\texttt{emergency\_\allowbreak exception} (boolean), (ii)~\texttt{emergency\_\allowbreak justification} (text summary of the identified risk), and (iii)~\texttt{post\_\allowbreak notice\_\allowbreak issued} (datetime). These records are available for regulatory inspection and judicial review. If the emergency exception is found to have been invoked without credible justification, the company faces aggravated liability --- the abuse of a safety mechanism to circumvent governance obligations is treated as a more serious breach than ordinary non-compliance.

\emph{Third, preservation of estate status.} Emergency deployment does not extinguish or downgrade the user's estate classification. A periodic tenant remains a periodic tenant, and a Stage~2 equitable interest holder retains that status. Their rights are temporarily overridden by the priority of life safety, not permanently diminished. Once the safety risk is resolved, the user retains the right to request rollback to a safe configuration or to receive compensation for irreversible changes.

Property rights do not override the right to life. This is not an innovation but a restatement of a principle as old as property law itself.}
